% problem-000235 6 醋酸亚铬无氧合成
\section*{6. 醋酸亚铬无氧合成}
醋酸亚铬水合物是一种存在金属-金属键的双核 $\mathrm{Cr(II)}$ 配合物，在潮湿或水溶液状态下很容易被空气中

的氧气氧化；而干燥的醋酸亚铬水合物固体被氧化的速率较慢。醋酸亚铬水合物可以通过 $Cr^{2+}$ 与醋酸根 $(\mathrm{AcO}^{-})$ 的配位直接合成。向 $CrCl_{3}$ 水溶液与 Zn 粒的混合物中加入浓盐酸，形成蓝色的 $CrCl_{2}$ 溶液；将该 $CrCl_{2}$ 溶液滴加至 HOAc-NaOAc 混合溶液中，加热形成深红色的醋酸亚铬配合物；醋酸亚铬在水中溶解度较低，冷却可使其从水溶液中析出。将析出的醋酸亚铬晶体通过减压过滤进行分离，用无氧的冰水和乙醚各洗涤固体 2 次，快速抽干即得到干燥的醋酸亚铬水合物晶体。

（图：）

\subsection*{6-1}
为防止氧化, 采用如右的仪器装置合成醋酸亚铬(只用于生成醋酸亚铬固体, 减压过滤不在该装置中进行)。其中 3 为氮气袋; 5 为玻璃砂芯隔板; A 为磨口玻璃接口, 可使 1 和 2 的组合装置整体 ${360}^{ \circ  }$ 旋转。该装置可在无氧条件下生成 ${\mathrm{{CrCl}}}_{2}$ 水溶液, 随后通过旋转接口 A, 将 ${\mathrm{{CrCl}}}_{2}$ 溶液转移至 HOAc-NaOAc 混合溶液中。

\subsection*{6-1}
-1 根据上述描述的步骤，写出实验开始时各试剂的位置，用图中的数字编号表示：

试剂 A: $\mathrm{CrCl}_{3}$ 水溶液与 $\mathrm{Zn}$ 粒的混合物; 位置编号: ( );

试剂 B: 浓盐酸; 位置编号: ( );

试剂 C: HOAc-NaOAc 混合溶液; 位置编号: ( )。

\subsection*{6-1}
-2 实验开始时, 先将体系抽真空, 去除其中的空气, 再充入 $\mathrm{N}_{2}$ , 随后向体系中加入浓盐酸, 此时须打开玻璃活塞 7。为避免空气扩散进入体系, 应采用何种措施? 6-1-3 玻璃砂芯隔板 5 的作用是什么? 6-1-4 醋酸亚铬水合物为抗磁性, 其中 Cr 均为六配位, 画出醋酸亚铬水合物的结构式 (本题画图时 AcO $^-$ 可用 $\circ$ 表示), 并指出 Cr-Cr 键的键级。 6-1-5 该水合物中 Cr-Cr $\pi$ 键来自于 Cr 的 d 轨道。假定 Cr-Cr 为 z 轴方向, 写出形成 $\pi$ 键的 d 轨道名称。

6-1-6 醋酸亚铬在水溶液中很容易被空气氧化, 其中一个产物为 $+ 1$ 价的三核 $\mathrm{Cr}^{3+}$ 配合物阳离子。该阳离子中$\mathrm{Cr}^{3+}$ 与 $\mathrm{AcO}^{-}$ 的比例与醋酸亚铬中 $\mathrm{Cr}$ 原子与 $\mathrm{AcO}^{-}$ 的比例相同, 三个 $\mathrm{Cr}^{3+}$ 环境等价, 均为 6 配位, 其中一个配体是水分子, 不存在 $\mathrm{Cr}-\mathrm{Cr}$ 键。画出该三核配合物的结构式。
6-2 无水醋酸亚铬可以通过醋酸亚铬水合物在真空中缓慢脱水获得, 但通常得到的是无定形的产物, 且纯度不高。将金属 $\mathrm{Cr}$ 粉末在冰醋酸中加热回流才能制备较纯的无水醋酸亚铬。
6-2-1 由于回流过程在开放体系中进行, 为有效避免环境中水汽的影响, 须向反应体系中添加一种物质。该物质是什么?
6-2-2 研究表明, 为加快反应速率, 可以向反应体系中添加少量的氢溴酸, 其作用是什么?
6-2-3 与醋酸亚铬水合物相比, 无水醋酸亚铬中 $\mathrm{Cr}-\mathrm{Cr}$ 键长 ( )。

（图：）

\subsection*{6-2}
-3 与醋酸亚铬水合物相比，无水醋酸亚铬中 Cr-Cr 键长（）。 (a) 更长 (b) 更短 (c) 相同
